\documentclass[11pt,a4paper]{article}
\usepackage[margin=1in]{geometry}
\usepackage{amsmath,amssymb,amsthm}
\usepackage{algorithm}
\usepackage{algpseudocode}
\usepackage{booktabs}
\usepackage{hyperref}

\theoremstyle{definition}
\newtheorem{definition}{Definition}[section]
\newtheorem{theorem}{Theorem}[section]
\newtheorem{lemma}[theorem]{Lemma}
\newtheorem{remark}{Remark}[section]

\title{\textbf{Theoretical Foundations, Skip Connection Dynamics, and Encoder-Decoder Synergy in U-Net Architecture}}
\author{Team Scaibu \\ \small Email: \texttt{jaydeep.9664920749@gmail.com} \\ \small Architecture and Core Engine Team}
\date{October 2026}

\begin{document}
\maketitle

\begin{abstract}
This document provides the formal mathematical specification, multi-resolution contracting-expanding paths, and long-range skip connection mechanics of the U-Net architecture. Dense semantic segmentation demands simultaneous high-level semantic context (global classification) and precise spatial pixel localization (boundary segmentation). U-Net achieves this through a contracting encoder pathway that downsamples spatial grids while increasing channel capacity, followed by a symmetric expanding decoder pathway that upsamples features and concatenates high-resolution encoder features via skip connections. We formulate the multi-scale recurrent operators, prove boundary reconstruction preservation, derive computational complexities, trace a worked numerical example, and document diffusion architecture usage.
\end{abstract}

\section{Notation and Mathematical Preliminaries}

Let $X \in \mathbb{R}^{C_{\text{in}} \times H \times W}$ denote the input image tensor where spatial height $H$ and width $W$ are divisible by $2^L$.

\begin{itemize}
    \item $L \in \mathbb{N}_{\ge 1}$: Total hierarchical depth levels ($L = 2$ in minimal specification).
    \item $C_1, C_2, C_{\text{out}} \in \mathbb{N}_{\ge 1}$: Channel dimensions at encoder level 1, bottleneck level 2, and output mask.
    \item $\mathcal{E}_l$: Contracting encoder operator at level $l$.
    \item $\mathcal{D}_l$: Expanding decoder operator at level $l$.
    \item $S_l \in \mathbb{R}^{C_l \times H_l \times W_l}$: High-resolution skip feature map cached from encoder level $l$.
    \item $\mathcal{U}_{2\times}(\cdot)$: Spatial upsampling operator (nearest-neighbor or transposed convolution).
    \item $\mathbin{\Vert}$: Channel-wise tensor concatenation operator.
    \item $Y \in \mathbb{R}^{C_{\text{out}} \times H \times W}$: Output dense segmentation probability or mask tensor.
\end{itemize}

\begin{table}[h]
\centering
\caption{Summary of U-Net Architecture Topology}
\begin{tabular}{lll}
\toprule
\textbf{Stage} & \textbf{Mathematical Transformation} & \textbf{Feature Map Shape} \\
\midrule
\textbf{Encoder Level 1} & $S_1 = \sigma(\mathcal{W}_{\text{enc1}} * X)$ & $(C_1, H, W)$ \\
\textbf{Downsampling} & $P_1 = \text{MaxPool}_{2 \times 2}(S_1)$ & $(C_1, H/2, W/2)$ \\
\textbf{Bottleneck (Level 2)} & $B = \sigma(\mathcal{W}_{\text{enc2}} * P_1)$ & $(C_2, H/2, W/2)$ \\
\textbf{Upsampling} & $U_1 = \mathcal{U}_{2 \times}(B)$ & $(C_2, H, W)$ \\
\textbf{Skip Fusion} & $F_1 = [U_1 \mathbin{\Vert} S_1]$ & $(C_2 + C_1, H, W)$ \\
\textbf{Decoder Level 1} & $Y = \sigma(\mathcal{W}_{\text{dec1}} * F_1)$ & $(C_{\text{out}}, H, W)$ \\
\bottomrule
\end{tabular}
\end{table}

\section{Problem Definition and Core Concepts}

\subsection{Intuitive Meaning}
Standard convolutional classifiers compress spatial resolution through successive pooling layers until spatial localization is destroyed. In image-to-image translation and biomedical image segmentation, the model must output predictions at the exact native input resolution. U-Net bridges the contracting encoder with the expanding decoder using direct identity skip connections, allowing sharp pixel-level boundary gradients to bypass the low-resolution bottleneck entirely.

\subsection{Formal Mathematical Recurrences}

\begin{align}
S_1 &= \sigma(\mathcal{W}_{\text{enc1}} * X) \\
P_1 &= \text{MaxPool}_{2 \times 2}(S_1) \\
B &= \sigma(\mathcal{W}_{\text{enc2}} * P_1) \\
U_1(c, 2h + v, 2w + u) &= B(c, h, w) \quad \forall v, u \in \{0, 1\} \\
F_1 &= [U_1 \mathbin{\Vert} S_1] \\
Y &= \sigma(\mathcal{W}_{\text{dec1}} * F_1)
\end{align}

\section{Algorithm Specification}

\begin{algorithm}[H]
\caption{U-Net Encoder-Decoder Forward Pass}
\label{alg:unet}
\begin{algorithmic}[1]
\Require Input $X \in \mathbb{R}^{C_{\text{in}} \times H \times W}$, weights $\mathcal{W}_{\text{enc1}}, \mathcal{W}_{\text{enc2}}, \mathcal{W}_{\text{dec1}}$.
\Ensure Dense segmentation tensor $Y \in \mathbb{R}^{C_{\text{out}} \times H \times W}$, skip $S_1$, bottleneck $B$.
\State $S_1 \gets \sigma(\text{Conv3x3}(X, \mathcal{W}_{\text{enc1}}, \text{pad}=1))$
\State $P_1 \gets \text{MaxPool2x2}(S_1)$
\State $B \gets \sigma(\text{Conv3x3}(P_1, \mathcal{W}_{\text{enc2}}, \text{pad}=1))$
\State $U_1 \gets \text{Upsample2xNearest}(B)$
\State $F_1 \gets \text{ConcatChannels}([U_1 \mathbin{\Vert} S_1])$
\State $Y \gets \sigma(\text{Conv3x3}(F_1, \mathcal{W}_{\text{dec1}}, \text{pad}=1))$
\State \Return $\{Y, S_1, B\}$
\end{algorithmic}
\end{algorithm}

\section{Theoretical Guarantees and Spatial Fidelity}

\begin{theorem}[Spatial Information Preservation via Long-Range Skips]
Let $f_{\text{enc}}$ and $f_{\text{dec}}$ denote the encoder and decoder pipelines. Under long-range skip connection concatenation $F_1 = [U_1 \mathbin{\Vert} S_1]$, high-frequency spatial gradients with wavevector $k > \pi / 2$ flow back to input layers with zero spatial downsampling attenuation:
\begin{equation}
\frac{\partial \mathcal{L}}{\partial X} \supseteq \frac{\partial \mathcal{L}}{\partial Y} \frac{\partial Y}{\partial S_1} \frac{\partial S_1}{\partial X}
\end{equation}
without traversing the sub-Nyquist bottleneck $P_1$.
\end{theorem}

\begin{proof}
Because $F_1 = [U_1 \mathbin{\Vert} S_1]$ separates channels into distinct sub-matrices, the Jacobian $\frac{\partial F_1}{\partial S_1} = [0 \mathbin{\Vert} I]^\top$ contains an exact identity block, allowing full-resolution gradients to propagate directly through $S_1$ without downsampling filtering.
\end{proof}

\subsection{Limitations}
\begin{itemize}
    \item \textbf{Input Resolution Divisibility}: Input spatial dimensions must be exact multiples of $2^L$ to ensure perfect spatial alignment during upsampling concatenation.
\end{itemize}

\section{Complexity Analysis}

\subsection{Time Complexity}
\begin{itemize}
    \item \textbf{Total Time Complexity}: $\Theta(H \cdot W \cdot (C_{\text{in}} C_1 + \frac{1}{4} C_1 C_2 + (C_2 + C_1) C_{\text{out}}))$ FLOPs.
\end{itemize}

\subsection{Space Complexity}
\begin{itemize}
    \item \textbf{Storage}: $\mathcal{O}(H \cdot W \cdot (C_1 + C_{\text{out}}) + \frac{1}{4} H W C_2)$ memory for activations and cached skip tensors.
\end{itemize}

\section{Exactness and Error Analysis}

Convolutions, max-pooling, and nearest-neighbor upsamplings are algebraically exact.

\section{Worked Numerical Example}

Let $C_{\text{in}} = 1, C_1 = 1, C_2 = 1, C_{\text{out}} = 1, H = 2, W = 2$.
Input $X = [[[1.0, 1.0], [1.0, 1.0]]]$.
Weights: all $1 \times 1$ filters with weight $[1.0]$.

\begin{enumerate}
    \item $S_1 = X = [[[1.0, 1.0], [1.0, 1.0]]]$.
    \item $P_1 = \text{MaxPool}(S_1) = [[[1.0]]] \in \mathbb{R}^{1 \times 1 \times 1}$.
    \item $B = [[[1.0]]]$.
    \item $U_1 = \text{Upsample2x}(B) = [[[1.0, 1.0], [1.0, 1.0]]]$.
    \item $F_1 = [U_1 \mathbin{\Vert} S_1] \in \mathbb{R}^{2 \times 2 \times 2}$ with channel 0 and 1 all ones.
    \item Decoder Conv (weight $[0.5, 0.5]$):
    $Y = 0.5(1.0) + 0.5(1.0) = 1.0 \implies Y = [[[1.0, 1.0], [1.0, 1.0]]]$.
\end{enumerate}

\section{Numerical Considerations and Edge Cases}

\begin{itemize}
    \item \textbf{Padding Modes}: Using zero-padding ($p=1$) in $3 \times 3$ convolutions guarantees that intermediate spatial dimensions match exact powers of 2 without border cropping.
\end{itemize}

\begin{thebibliography}{9}
\bibitem{ronneberger2015u} O.~Ronneberger, P.~Fischer, and T.~Brox, ``U-Net: Convolutional Networks for Biomedical Image Segmentation,'' in \emph{Medical Image Computing and Computer-Assisted Intervention (MICCAI)}, pp.~234--241, 2015.
\bibitem{ho2020denoising} J.~Ho, A.~Jain, and P.~Abbeel, ``Denoising Diffusion Probabilistic Models,'' in \emph{Advances in Neural Information Processing Systems (NeurIPS)}, pp.~6840--6851, 2020.
\end{thebibliography}

\end{document}
